\VerdierPhysicalPageAnchor{0051}%
\VerdierSemanticAnchor{ch1-1-5}
\subsection{\texorpdfstring{$\mathbb Z^n$}{Z n}-catégories}

\VerdierSemanticAnchor{ch1-1-5-1}
\noindent\textbf{1.5.1.}\quad
Soient $(\mathcal C,\mathbb Z^n,T,c)$ une $\mathbb Z^n$-catégorie,
$(\mathcal C',\mathbb Z^n,T',c')$ une $\mathbb Z^n$-catégorie
additive,
\[
  (F,p)\colon(\mathcal C,\mathbb Z^n,T,c)
  \longrightarrow(\mathcal C',\mathbb Z^n,T',c')
\]
un $\mathbb Z^n$-foncteur tordu par un cocycle normalisé
$\alpha\in Z^2(\mathbb Z^n,\mathbb Z/2\mathbb Z)$. Désignons par
$e_1,\ldots,e_n$ la base canonique de $\mathbb Z^n$. On a donc pour
tout $i\in[1,n]$ un isomorphisme de foncteurs :

\VerdierSemanticAnchor{ch1-1-5-1-eq1}
\begin{equation}
  p(e_i)\colon FT(e_i)\longrightarrow T'(e_i)F
  \tag{1.5.1.1}
\end{equation}
et pour tout couple $(i,j)$, $i\ne j$, un diagramme
$\alpha(e_i,e_j)\alpha^{-1}(e_j,e_i)$-commutatif :

\VerdierSemanticAnchor{ch1-1-5-1-eq2}
\begin{center}
  \resizebox{0.98\textwidth}{!}{%
    $\begin{tikzcd}[ampersand replacement=\&,column sep=4.2em,row sep=3.3em]
      FT(e_i+e_j)
        \arrow[r,"{F*c(e_i,e_j)}"]
        \arrow[d,"{F*c(e_j,e_i)}"]
        \&
      FT(e_j)T(e_i)
        \arrow[r,"{p(e_j)*T(e_i)}"]
        \&
      T'(e_j)FT(e_i)
        \arrow[d,"{T'(e_j)*p(e_i)}"'] \\
      FT(e_i)T(e_j)
        \arrow[d,"{p(e_i)*T(e_j)}"]
        \& {} \&
      T'(e_j)T'(e_i)F
        \arrow[d,"{c'(e_i,e_j)^{-1}*F}"'] \\
      T'(e_i)FT(e_j)
        \arrow[r,"{T'(e_i)*p(e_j)}"']
        \&
      T'(e_i)T'(e_j)F
        \arrow[r,"{c'(e_j,e_i)^{-1}*F}"']
        \&
      T'(e_i+e_j)F
    \end{tikzcd}$%
  }\;.
  \hfill\textup{(1.5.1.2)}
\end{center}

\VerdierSemanticAnchor{ch1-1-5-2}
\noindent\textbf{Proposition 1.5.2.}\quad
\emph{Soient $F\colon\mathcal C\to\mathcal C'$ un foncteur et pour
tout $i\in[1,n]$ un isomorphisme de foncteurs :
\[
  p_i\colon FT(e_i)\longrightarrow T'(e_i)F
\]
tel que pour tout couple $i,j$, $i\ne j$, le diagramme analogue au
diagramme (1.5.1.2), où l'on remplace les $p(e_i)$ par les $p_i$, soit
$\alpha(e_i,e_j)\alpha^{-1}(e_j,e_i)$-commutatif. Il existe un et un
seul $\mathbb Z^n$-foncteur tordu par le cocycle $\alpha$ :
\[
  (F,p)\colon(\mathcal C,\mathbb Z^n,T,c)
  \longrightarrow(\mathcal C',\mathbb Z^n,T',c')
\]
tel que :
\[
  p(e_i)=p_i,
  \qquad i\in[1,n].
\]}

\VerdierSemanticAnchor{ch1-1-5-2-proof-intro}
La démonstration de l'unicité se fait en examinant les diagrammes
(1.4.3.11). Nous laissons au lecteur le soin de vérifier ce point. Nous
nous bornerons à donner quelques indications sur la démonstration de
l'existence. Celle-ci se fait par récurrence sur l'entier $n$.
